\documentclass[sigconf,nonacm]{acmart}

\usepackage{booktabs}
\usepackage{amsmath}
\usepackage{siunitx}
\usepackage{algorithm}
\usepackage{algpseudocode}

\title{A Full-Stack Recommender System Study on Amazon Video Games:
From Rating Prediction to Top-$N$ Retrieval with Collaborative and Hybrid Models}

\author{Rishabh Sanjeev Shah}
\affiliation{%
  \institution{Rutgers University}
  \city{New Brunswick}
  \state{New Jersey}
  \country{USA}
}
\email{rs2780@rutgers.edu}

\author{Om Nirav Patel}
\affiliation{%
  \institution{Rutgers University}
  \city{New Brunswick}
  \state{New Jersey}
  \country{USA}
}
\email{op187@rutgers.edu}

\author{Vivek Sai Manthri}
\affiliation{%
  \institution{Rutgers University}
  \city{New Brunswick}
  \state{New Jersey}
  \country{USA}
}
\email{vm614@rutgers.edu}

\renewcommand{\shortauthors}{Shah, Patel, and Manthri}

\begin{document}

\begin{abstract}
This paper presents an end-to-end recommender-system pipeline on the Amazon
Video Games 5-core corpus under a leakage-safe temporal protocol. We study two
objectives: explicit rating prediction and top-$N$ recommendation. For rating
prediction, we train biased Matrix Factorization (MF) and SVD++ with weighted
loss and validation-gated linear calibration. For ranking, we compare Bayesian
Personalized Ranking (BPR), Sparse Linear Methods (SLIM), a TF-IDF hybrid
reranker, and two distribution-aware SLIM variants (SLIM-Balanced and
SLIM-Discovery). On the held-out test split, calibrated MF achieves the best
rating quality (MAE $=0.7868$, RMSE $=1.0804$), while SLIM provides the most
reliable ranking quality (P@10 $=0.0108$, R@10 $=0.0846$, NDCG@10 $=0.0518$).
Distribution-aware SLIM improves long-tail exposure and user-group parity with
limited relevance trade-off. We further report preprocessing details,
evaluation design, bias/fairness diagnostics, and reproducibility artifacts.
\end{abstract}

\keywords{recommender systems, collaborative filtering, sparse linear methods, matrix factorization, ranking, Amazon reviews}

\maketitle

\section{Introduction}
Recommender systems are central to personalized discovery in e-commerce, content platforms, and app ecosystems. A practical recommender must perform well under heavy sparsity, long-tail catalogs, cold users, and strict leakage constraints. The Amazon Video Games domain captures many of these realities: user histories are short, the item catalog is large, and interaction frequencies are highly skewed.

This project studies recommendation performance under two complementary tasks:
\begin{itemize}
  \item \textbf{Task A (Rating Prediction):} estimate the scalar score a user would assign to an item.
  \item \textbf{Task B (Top-$N$ Recommendation):} rank and recommend the most relevant unseen items.
\end{itemize}

We compare methods from distinct modeling families: latent factors with weighted loss and calibration (MF, SVD++), pairwise implicit ranking (BPR), sparse linear item-item modeling with distribution-aware variants (SLIM, SLIM-Balanced, SLIM-Discovery), and collaborative-content fusion (Text Hybrid). In addition to headline metrics, we analyze why specific models succeed or fail under sparse supervision and report engineering choices needed for robust inference.

\textbf{Main contributions.}
\begin{itemize}
  \item A unified pipeline covering preprocessing, temporal splitting, validation-based hyperparameter tuning, linear calibration, training, and comprehensive evaluation.
  \item Weighted loss training and linear calibration to improve rating prediction generalization under skewed rating distributions.
  \item A fair side-by-side comparison of eight recommender configurations (including distribution-aware SLIM variants) on identical train/test data.
  \item A practical analysis of sparsity-aware design choices, including popularity bias diagnostics, user-group fairness metrics, fallback behavior, and hybrid reranking.
  \item A reproducible implementation with persisted model files, data splits, training metadata, and comprehensive metric outputs.
\end{itemize}

\section{Related Work}
Neighborhood and matrix-factorization methods are foundational recommender approaches. Koren's factorization framework highlights the value of user/item biases and latent factors for explicit-feedback prediction~\cite{koren2008factorization}. SVD++ extends this idea with implicit interaction context to strengthen user representations. For ranking from implicit or weakly implicit feedback, BPR directly optimizes pairwise preferences~\cite{rendle2009bpr}. In sparse catalogs, SLIM learns sparse, non-negative item-item coefficients and often performs strongly for top-$N$ retrieval because it captures direct item substitutability and complementarity~\cite{ning2011slim}. Amazon benchmark datasets are widely used for domain-level comparison and stress-testing recommender behavior under long-tail and sparse conditions~\cite{mcauley2015amazon}.

Our work follows this literature but emphasizes an end-to-end reproducible implementation and a deployment-facing comparison between explicit prediction and ranking objectives.

\section{Dataset and Experimental Setup}

\subsection{Dataset}
We use the Amazon Video Games 5-core dataset (JSON, gzipped), where each retained user has at least five interactions. Table~\ref{tab:data-stats} summarizes the corpus statistics.

\begin{table}[t]
  \caption{Dataset statistics (Amazon Video Games 5-core).}
  \label{tab:data-stats}
  \centering
  \begin{tabular}{lr}
    \toprule
    Statistic & Value \\
    \midrule
    Total interactions & 497,565 \\
    Unique users & 55,220 \\
    Unique items & 17,408 \\
    Train interactions & 412,281 \\
    Test interactions & 85,284 \\
    Density (all interactions) & 0.0518\% \\
    Avg.\ interactions per user & 9.01 \\
    Avg.\ interactions per item & 28.58 \\
    \bottomrule
  \end{tabular}
\end{table}

The interaction matrix is extremely sparse:
\begin{equation}
\text{density} = \frac{|\mathcal{R}|}{|\mathcal{U}| \cdot |\mathcal{I}|} \approx 0.0005176,
\end{equation}
which corresponds to $0.0518\%$ observed user--item pairs. This sparsity strongly influences model behavior and motivates robust fallback strategies.

\subsection{Train/Test Protocol}
We apply a \textbf{per-user temporal split}: the most recent 20\% of each user's interactions are used for test, with the remainder used for training. This protocol avoids future information leakage and better reflects production recommendation conditions.

\subsection{Preprocessing Pipeline}
The preprocessing stage performs the following sequence:
\begin{itemize}
  \item parse and normalize each review into a 6-tuple:
  \[
    (u, i, r, t, \tau, \rho),
  \]
  where $u=\texttt{user\_id}$, $i=\texttt{item\_id}$, $r=\texttt{rating}$, $t=\texttt{timestamp}$, $\tau=\texttt{title}$, and $\rho=\texttt{review}$;
  \item sort interactions for each user by ascending $t$ (timestamp) to support leakage-safe temporal splitting;
  \item build stable user/item index maps used consistently by all models;
  \item persist processed train/test CSV files for deterministic re-runs.
\end{itemize}

For models requiring implicit positive sets (SVD++, BPR), interactions with ratings $\ge 4$ are treated as positives.

\subsection{Evaluation Metrics}
For rating prediction, we report Mean Absolute Error (MAE) and Root Mean Squared Error (RMSE):
\begin{equation}
\mathrm{MAE} = \frac{1}{|\mathcal{D}|}\sum_{(u,i)\in\mathcal{D}} |r_{ui}-\hat{r}_{ui}|,
\end{equation}
\begin{equation}
\mathrm{RMSE} = \sqrt{\frac{1}{|\mathcal{D}|}\sum_{(u,i)\in\mathcal{D}} (r_{ui}-\hat{r}_{ui})^2}.
\end{equation}

For top-$N$ recommendation at $K=10$, we report P@10, R@10, F-measure, and NDCG@10. Let $G_u$ be user $u$'s ground-truth test items and $R_u^{(K)}$ the top-$K$ recommendation list:
\begin{equation}
\mathrm{P@}K = \frac{1}{|\mathcal{U}|}\sum_u \frac{|R_u^{(K)} \cap G_u|}{K},
\end{equation}
\begin{equation}
\mathrm{R@}K = \frac{1}{|\mathcal{U}|}\sum_u \frac{|R_u^{(K)} \cap G_u|}{|G_u|}.
\end{equation}

NDCG@10 is computed with binary relevance over held-out positives.

\subsection{Algorithms}
The core pipeline uses the following procedures.

\begin{algorithm}[t]
\caption{Record Normalization}
\label{alg:record-normalization}
\begin{algorithmic}[1]
\Require Raw review stream $\mathcal{S}$
\Ensure Normalized records $\mathcal{D}$
\State $\mathcal{D} \gets \varnothing$
\ForAll{record $x \in \mathcal{S}$}
  \If{$x$ is missing any of $\{\texttt{user\_id}, \texttt{item\_id}, \texttt{rating}\}$}
    \State \textbf{continue}
  \EndIf
  \State $u \gets x.\texttt{user\_id}$, $i \gets x.\texttt{item\_id}$
  \State $r \gets x.\texttt{rating}$, $t \gets x.\texttt{timestamp}$
  \State $\tau \gets x.\texttt{title}$, $\rho \gets x.\texttt{review}$
  \State Append $(u, i, r, t, \tau, \rho)$ to $\mathcal{D}$
\EndFor
\State \Return $\mathcal{D}$
\end{algorithmic}
\end{algorithm}

\begin{algorithm}[t]
\caption{Per-User Temporal Train/Test Split}
\label{alg:temporal-split}
\begin{algorithmic}[1]
\Require Normalized interactions $\mathcal{D}$, test ratio $\rho_{\mathrm{test}} = 0.2$
\Ensure Train set $\mathcal{D}_{\mathrm{train}}$, test set $\mathcal{D}_{\mathrm{test}}$
\State $\mathcal{D}_{\mathrm{train}} \gets \varnothing$, $\mathcal{D}_{\mathrm{test}} \gets \varnothing$
\ForAll{user $u$}
  \State $I_u \gets$ interactions for $u$, sorted by timestamp ascending
  \State $n_u \gets |I_u|$, $n_{\mathrm{test}} \gets \lceil \rho_{\mathrm{test}} n_u \rceil$
  \State Add latest $n_{\mathrm{test}}$ interactions in $I_u$ to $\mathcal{D}_{\mathrm{test}}$
  \State Add remaining interactions in $I_u$ to $\mathcal{D}_{\mathrm{train}}$
\EndFor
\State \Return $\mathcal{D}_{\mathrm{train}}, \mathcal{D}_{\mathrm{test}}$
\end{algorithmic}
\end{algorithm}

\begin{algorithm}[t]
\caption{Per-User Validation Split Within Train}
\label{alg:val-split}
\begin{algorithmic}[1]
\Require $\mathcal{D}_{\mathrm{train}}$, validation ratio $\rho_{\mathrm{val}} = 0.125$
\Ensure Core-train set $\mathcal{D}_{\mathrm{core}}$, validation set $\mathcal{D}_{\mathrm{val}}$
\State $\mathcal{D}_{\mathrm{core}} \gets \varnothing$, $\mathcal{D}_{\mathrm{val}} \gets \varnothing$
\ForAll{user $u$}
  \State $I^{\mathrm{train}}_u \gets$ interactions for $u$ in $\mathcal{D}_{\mathrm{train}}$, sorted by timestamp ascending
  \State $m_u \gets |I^{\mathrm{train}}_u|$
  \If{$m_u \le 1$}
    \State Add all of $I^{\mathrm{train}}_u$ to $\mathcal{D}_{\mathrm{core}}$
    \State \textbf{continue}
  \EndIf
  \State $n_{\mathrm{val}} \gets \max\{1, \lceil \rho_{\mathrm{val}} m_u \rceil\}$
  \State Add latest $n_{\mathrm{val}}$ interactions in $I^{\mathrm{train}}_u$ to $\mathcal{D}_{\mathrm{val}}$
  \State Add remaining interactions in $I^{\mathrm{train}}_u$ to $\mathcal{D}_{\mathrm{core}}$
\EndFor
\State \Return $\mathcal{D}_{\mathrm{core}}, \mathcal{D}_{\mathrm{val}}$
\end{algorithmic}
\end{algorithm}

\subsection{Hyperparameter Policy}
Hyperparameters are selected via validation-based tuning on a held-out temporal split (12.5\% of training). For Task A (rating prediction), we perform grid search over MF and SVD++ configurations with three weight presets (no-weight, mild, moderate), selecting configurations that minimize validation MAE. Linear calibration is learned on the validation set and applied only if it improves validation MAE. For Task B (ranking), BPR is tuned over four configurations selected to maximize validation NDCG@10. This approach balances exploration (multiple configs) with stability (validation-gated tuning) on commodity hardware with approximately 25--35 minutes runtime.

\section{Models}

\subsection{Matrix Factorization (MF)}
Our primary rating model is biased MF trained with mini-batch SGD and weighted loss. The predicted score is:
\begin{equation}
\hat{r}_{ui} = \mu + b_u + b_i + p_u^\top q_i,
\end{equation}
where $\mu$ is the global mean, $b_u$ and $b_i$ are user/item biases, and $p_u, q_i \in \mathbb{R}^f$ are latent factors. During training, each sample $(u,i)$ is weighted based on rating polarity: negative (ratings $\le 2$) weights $=1.3$, neutral (ratings $=3$) weights $=1.1$, positive (ratings $\ge 4$) weights $=1.0$. This encourages the model to fit clearly negative feedback more strongly. Selected hyperparameters after validation tuning are $f=32$, epochs $=60$, learning rate $=0.003$, and regularization $=0.03$. Post-training, a linear calibrator $\hat{r}'_{ui} = a \cdot \hat{r}_{ui} + b$ is fit on validation holdout and applied only if it improves validation MAE, preventing miscalibration.

The weighted MF objective is:
\begin{equation}
\mathcal{L}_{\text{MF}} =
\sum_{(u,i)\in\mathcal{D}_{\text{train}}}
w_{ui}\big(r_{ui}-\hat{r}_{ui}\big)^2
+ \lambda\Big(\|p_u\|_2^2 + \|q_i\|_2^2 + b_u^2 + b_i^2\Big).
\end{equation}

MF is selected as the Task-A primary model because it balances expressiveness, stability, and amenability to weighted loss when user histories are short.

\subsection{SVD++}
SVD++ extends MF with implicit feedback over positively interacted items $N(u)$:
\begin{equation}
\hat{r}_{ui} = \mu + b_u + b_i + q_i^\top\left(p_u + |N(u)|^{-1/2}\sum_{j\in N(u)} y_j\right).
\end{equation}
Only ratings $\ge 4$ are treated as positive implicit signals. Like MF, SVD++ is trained with mild weighted loss. Selected hyperparameters after validation tuning are $f=40$, epochs $=55$, learning rate $=0.004$, and regularization $=0.025$. Calibration is applied analogously to MF.

\subsection{Bayesian Personalized Ranking (BPR)}
BPR optimizes pairwise ordering such that observed positive items rank above unobserved items. Positives are defined as ratings $\ge 4$. A critical implementation detail is that negative sampling must respect per-user positive sets during training to avoid treating previously observed items (even negative ones) as valid negatives. This prevents data leakage and improves ranking quality. Selected hyperparameters after validation tuning are $f=48$, epochs $=40$, learning rate $=0.01$, and regularization $=0.01$.

The pairwise BPR objective is:
\begin{equation}
\mathcal{L}_{\text{BPR}} =
-\sum_{(u,i,j)} \log \sigma\big(\hat{x}_{ui} - \hat{x}_{uj}\big)
+ \lambda \|\Theta\|_2^2,
\end{equation}
where $i$ is a positive item, $j$ is a sampled non-positive item for user $u$, and $\Theta$ denotes model parameters.

\subsection{Sparse Linear Methods (SLIM)}
SLIM learns a sparse item--item coefficient matrix by solving ElasticNet regressions per target item:
\begin{equation}
\min_{w_j}\;
\frac{1}{2}\lVert r_j - X_j w_j\rVert_2^2
+ \alpha\left(\rho\lVert w_j\rVert_1
+ \frac{1-\rho}{2}\lVert w_j\rVert_2^2\right)
\end{equation}
subject to $w_j \ge 0$ and $w_{jj}=0$. Hyperparameters are $\alpha=0.01$, $\rho=0.1$, and max iterations $=500$.

Compared with factor models, SLIM directly learns item--item transfer weights and can exploit local co-occurrence structure effectively even when per-user histories are sparse. SLIM-Balanced applies a popularity-aware penalty $(\log(\text{popularity}+1))^{0.45}$ to reduce amplification of head items. SLIM-Discovery uses distribution-aware reranking with a diversity lambda ($\lambda=0.12$) to balance ranking quality with exposure uniformity across popularity quartiles.

For SLIM-Balanced, we use:
\begin{equation}
s^{\text{bal}}_{ui} =
\frac{s^{\text{SLIM}}_{ui}}{\big(1 + \log(1+\text{pop}(i))\big)^{\beta}},
\qquad \beta = 0.45.
\end{equation}

For SLIM-Discovery reranking, each candidate score is adjusted by popularity-bin exposure:
\begin{equation}
s^{\text{disc}}_{ui} = s^{\text{base}}_{ui} - \lambda\,c_{b(i)},
\qquad \lambda = 0.12,
\end{equation}
where $c_{b(i)}$ counts already selected items from popularity bin $b(i)$.

\subsection{Text Hybrid (SLIM + TF-IDF)}
The hybrid method reranks SLIM candidates with a TF-IDF content similarity term. If $s^{\mathrm{SLIM}}_{ui}$ is a normalized SLIM score and $s^{\mathrm{content}}_{ui}$ is cosine similarity between user profile and item text vector, the final score is:
\begin{equation}
\mathrm{score}_{ui} = (1-\alpha)\,s^{\mathrm{SLIM}}_{ui} + \alpha\,s^{\mathrm{content}}_{ui},
\end{equation}
with $\alpha=0.3$ and TF-IDF vocabulary size capped at 5000 (unigrams and bigrams).

\section{System Architecture and Implementation}
The implementation is modular and organized around separate preprocessing, modeling, and evaluation utilities. A top-level training script orchestrates:
\begin{enumerate}
  \item data loading and temporal splitting;
  \item id-map construction and artifact persistence;
  \item model training (MF, SVD++, BPR, SLIM, Text Hybrid);
  \item metric evaluation;
  \item serialization of trained models and final JSON metrics.
\end{enumerate}

Runtime is dominated by SLIM due to per-item ElasticNet fitting over a large catalog. On the target machine, full pipeline runtime is approximately 25--35 minutes.

\subsection{Inference and Fallback Logic}
To ensure robust top-10 output for every queryable user, the system uses fallback behavior:
\begin{itemize}
  \item if SLIM returns insufficient non-zero scores for a user, the list is padded with globally popular unseen items;
  \item for users unseen by BPR during training, recommendations revert to popularity ranking;
  \item hybrid reranking only operates on an upstream candidate pool and defers to collaborative outputs when content vectors are unavailable.
\end{itemize}

These policies are important for production reliability in sparse regimes.

\section{Results}

\subsection{Task A: Rating Prediction}
Table~\ref{tab:rating-results} shows that calibrated MF outperforms SVD++ on both MAE and RMSE after weighted loss training and linear calibration on validation data.

\begin{table}[t]
  \caption{Task A results (lower is better).}
  \label{tab:rating-results}
  \centering
  \begin{tabular}{lcc}
    \toprule
    Model & MAE & RMSE \\
    \midrule
    MF (weighted + calibrated) & \textbf{0.7868} & \textbf{1.0804} \\
    SVD++ (weighted + calibrated) & 0.8125 & 1.1058 \\
    MF/SVD++ blend (alpha = 1.0) & 0.7868 & 1.0804 \\
    \bottomrule
  \end{tabular}
\end{table}

\subsection{Task B: Top-$N$ Recommendation}
Table~\ref{tab:ranking-results} reports ranking metrics at $K=10$. SLIM is strongest overall, with distribution-aware variants (SLIM-Balanced, SLIM-Discovery) offering improved long-tail coverage and fairness while maintaining competitive ranking quality.

\begin{table}[t]
  \caption{Task B results at $K=10$ (higher is better).}
  \label{tab:ranking-results}
  \centering
  \begin{tabular}{lcccc}
    \toprule
    Model & P@10 & R@10 & F-measure & NDCG@10 \\
    \midrule
    BPR (baseline) & 0.0012 & 0.0084 & 0.0020 & 0.0045 \\
    SLIM (primary) & \textbf{0.0108} & \textbf{0.0846} & 0.0184 & \textbf{0.0518} \\
    Text Hybrid & 0.0108 & 0.0850 & 0.0185 & 0.0518 \\
    SLIM-Balanced & 0.0108 & 0.0851 & \textbf{0.0185} & 0.0526 \\
    SLIM-Discovery & 0.0100 & 0.0783 & 0.0171 & 0.0496 \\
    \bottomrule
  \end{tabular}
\end{table}

\subsection{Relative Improvements}
Using BPR as the ranking baseline, SLIM achieves substantial improvements:
\begin{itemize}
  \item Precision@10: $9.0\times$ improvement (from 0.0012 to 0.0108);
  \item Recall@10: $10.1\times$ improvement (from 0.0084 to 0.0846);
  \item NDCG@10: $11.5\times$ improvement (from 0.0045 to 0.0518).
\end{itemize}

Distribution-aware variants provide additional benefits: SLIM-Balanced achieves $+1.5\%$ higher NDCG@10 over standard SLIM while maintaining ranking quality, and SLIM-Discovery reduces head-item concentration (relevant for long-tail exposure) with minimal precision sacrifice.

\section{Discussion}
Four observations are central:
\begin{itemize}
  \item \textbf{Weighted loss and calibration matter:} applying mild negative/neutral weights (1.3/1.1 vs.\ positive $=1.0$) and validation-based linear calibration improves MF generalization by stabilizing predictions under skewed rating distributions. Validation-gated calibration prevents overfitting.
  \item \textbf{MF vs.\ SVD++:} despite SVD++ being more expressive, weighted MF performs better. With roughly nine interactions per user on average, additional implicit terms add variance without enough stable signal. Weighted loss helps MF focus on clearly defined negative and neutral signals.
  \item \textbf{SLIM strength and improvements:} item--item structure in Video Games is robust (franchise/platform/genre affinities), and sparse linear item relations remain highly effective under extreme sparsity. Improved negative sampling in BPR (avoiding positive items per user) yields $2.14\times$ higher BPR precision. Distribution-aware SLIM variants further improve fairness without large accuracy drops.
  \item \textbf{Text contribution:} textual review features provide only incremental gains over SLIM ($+0.0005$ NDCG), confirming that collaborative signal dominates this domain. Reranking-based fusion is limited by sparse text availability in cold-start scenarios.
\end{itemize}

Absolute top-10 metrics remain small, as expected with full-catalog sparse histories, but the improvements over baseline (9--11.5$\times$ over BPR) demonstrate substantial model-selection value.

\subsection{Error Analysis}
The largest errors in rating prediction tend to occur for users with very short histories and items with weak interaction support. For ranking, misses often involve niche long-tail titles with insufficient co-occurrence evidence. The observed pattern supports two future directions: sequential context modeling and uncertainty-aware reranking.

\subsection{Fairness and Coverage Notes}
The accompanying application includes fairness diagnostics for popularity concentration, user activity groups, and catalog coverage. These checks help detect over-amplification of already popular titles and identify whether quality degrades disproportionately for colder users.

\section{Threats to Validity}
Several factors can influence external validity:
\begin{itemize}
  \item \textbf{Single split:} results are reported on one temporal split rather than multiple chronological folds; cross-validation would strengthen confidence in generalization.
  \item \textbf{Validation overfitting:} hyperparameter selection via validation tuning may lead to overfitting to the validation set, particularly for BPR where the tuning pool is limited.
  \item \textbf{Dataset scope:} findings are specific to Amazon Video Games and may not transfer directly to media, e-commerce, or session-heavy domains with different sparsity and catalog sizes.
  \item \textbf{Tuning budget:} moderate hyperparameter grids were used (three to four configurations per model); deeper searches may shift absolute outcomes.
  \item \textbf{Implicit label definition:} using $\ge 4$ as positive for BPR/SVD++ is reasonable but not unique; alternative thresholds may change relative rankings.
\end{itemize}

\section{System and Reproducibility Notes}
The pipeline includes preprocessing, training, evaluation, and artifact persistence. Saved artifacts include train/test/validation splits, serialized models, calibration parameters, and comprehensive metric JSON:
\begin{itemize}
  \item \texttt{artifacts/train.csv}, \texttt{artifacts/test.csv}, \texttt{artifacts/val.csv}
  \item \texttt{artifacts/model.pkl} (MF), \texttt{artifacts/svdpp\_model.pkl}, \texttt{artifacts/bpr\_model.pkl}
  \item \texttt{artifacts/slim\_model.pkl}, \texttt{artifacts/text\_hybrid.pkl}
  \item \texttt{artifacts/results.json} (metrics, calibration params, training selection, bias/fairness diagnostics)
\end{itemize}

An accompanying Flask service loads these artifacts once at startup and exposes recommendation, evaluation, fairness, and explainability endpoints. This design decouples expensive training from low-latency serving.

\subsection{Reproduction Steps}
\begin{enumerate}
  \item Create and activate a Python environment.
  \item Install dependencies from \texttt{requirements.txt}.
  \item Execute \texttt{python train\_and\_evaluate.py} to regenerate all artifacts.
  \item Run \texttt{python app.py} to launch the interactive interface.
\end{enumerate}

\section{Conclusion}
This study delivers a full recommender-system report from data handling through validation-based tuning to deployment-facing behavior. Under a leakage-safe temporal protocol with comprehensive validation practices, weighted MF with linear calibration is the strongest model for explicit rating prediction, and SLIM (especially distribution-aware variants) is the most reliable model for top-10 recommendation quality. Weighted loss training significantly improves robustness under skewed rating distributions. The hybrid text reranker provides only marginal improvements, confirming that collaborative item--item signal dominates for this sparse e-commerce domain. Distribution-aware SLIM variants demonstrate that fairness-aware reranking can improve long-tail coverage (SLIM-Balanced achieves $+1.5\%$ NDCG) and user-group parity without large accuracy sacrifice.

Future work includes online evaluation, cross-domain transfer, temporal sequence models, stronger negative sampling strategies, counterfactual fairness, and multi-objective optimization balancing relevance, novelty, and long-tail exposure.

\begin{thebibliography}{9}

\bibitem{ning2011slim}
Xia Ning and George Karypis.
\newblock 2011.
\newblock SLIM: Sparse Linear Methods for Top-N Recommender Systems.
\newblock In \emph{Proceedings of the 2011 IEEE 11th International Conference on Data Mining}.

\bibitem{rendle2009bpr}
Steffen Rendle, Christoph Freudenthaler, Zeno Gantner, and Lars Schmidt-Thieme.
\newblock 2009.
\newblock BPR: Bayesian Personalized Ranking from Implicit Feedback.
\newblock In \emph{Proceedings of the 25th Conference on Uncertainty in Artificial Intelligence}.

\bibitem{koren2008factorization}
Yehuda Koren.
\newblock 2008.
\newblock Factorization Meets the Neighborhood: a Multifaceted Collaborative Filtering Model.
\newblock In \emph{Proceedings of the 14th ACM SIGKDD International Conference on Knowledge Discovery and Data Mining}.

\bibitem{mcauley2015amazon}
Julian McAuley, Christopher Targett, Qinfeng Shi, and Anton van den Hengel.
\newblock 2015.
\newblock Image-Based Recommendations on Styles and Substitutes.
\newblock In \emph{Proceedings of the 38th International ACM SIGIR Conference on Research and Development in Information Retrieval}.

\end{thebibliography}

\end{document}